% !TeX root = ../main.tex
\section{Background and Related Work}
\label{sec:related}

\paragraph{MI and simulated clients.}
Because MI sessions are confidential and expert coding is slow, work on computational MI
has increasingly used LLM-simulated clients as both a training environment and an
evaluation harness. \citet{yosef2024assessing} showed that an LLM prompted with the
session-satisfaction and working-alliance instruments used here produces ratings that are
statistically reliable and that separate three levels of therapist expertise --- the
validation basis on which those two instruments are used as an LLM evaluator, and the
reason we treat them as a satisfaction/alliance proxy rather than as clinical outcome
measures. \todo{cite 1--2 more recent MI-simulation papers (2025--2026) so this does not
read as lab-internal; check \citet{steenstra2024scaffolding}-style work and MI-coding LLM
papers.}

\paragraph{Preference optimization from self-generated data.}
DPO \citep{rafailov2023dpo} removed the explicit reward model from RLHF
\citep{christiano2017deeprl, ouyang2022instructgpt} and made it practical to train on
preference pairs produced by the model itself and labelled by an automatic judge
\citep{yuan2024selfrewarding, guo2024oaif, pace2024westofn}. Where the pairs come from
matters as much as the loss: search-based construction --- MCTS over reasoning traces
\citep{xie2024mctsdpo}, preference trees \citep{yuan2024eurus}, and dialogue-policy
planning \citep{yu2023promptmcts, chen2025broaden} --- consistently beats i.i.d.\ sampling
in domains where the payoff is delayed. PTO \citep{baruch2025pto} is the MI instance of
this idea: it grows a preference tree over therapist turns and scores each branch after
$K$ further simulated turns, so a candidate is judged by where it \emph{leads} rather than
by how it \emph{looks}. That work established PTO against an untrained baseline; it did not
compare it to an on-policy RL alternative, which is what we do here.

\paragraph{Group-relative optimization.}
GRPO \citep{shao2024deepseekmath} replaces the learned value function of PPO
\citep{schulman2017ppo} with a group baseline: sample $G$ completions for a prompt, and
weight each one's gradient by its reward standardized within the group. It has become the
default recipe for verifiable-reward domains, which makes ``does it transfer to a soft,
LLM-judged, multi-turn domain?'' the natural question --- and one that has to be asked at a
\emph{matched} candidate budget, since PTO's $M$ branches and GRPO's $G$ completions are
the same expenditure of generation and grading.

\paragraph{Reward hacking and LLM judges.}
Optimizing against a learned or prompted evaluator invites the policy to exploit the
evaluator rather than the task \citep{gao2023scaling, skalse2022defining}. In the RLHF
setting the best-documented instance is sycophancy \citep{sharma2024sycophancy,
perez2023discovering}, and LLM judges are known to prefer longer, more assertive and
more flattering answers \citep{zheng2023judging, dubois2024lengthcontrolled}. MI is an
unusually sharp test case for this, because the failure mode has a \emph{clinical} name:
an MI session in which the counsellor praises the client and stops asking questions is not
merely a worse answer, it is MI-inconsistent practice as the MITI manual
\citep{moyers2016miti} defines it. Our contribution here is less the observation that this
happens than the measurement: we quantify it on an instrument outside the reward
(\S\ref{sec:hacking}) and, more sharply, as a train/test generalization ratio across
graders (\S\ref{sec:validity}).

\paragraph{Judge reliability.}
Single-judge evaluation is now routinely criticised for exactly the coupling we have here
--- the same model family generates and grades \citep{panickssery2024selfpreference}. The
standard mitigations are multiple judges and human spot-checks. We report the reliability
of both graders explicitly (ICC over repeated scorings) so that cross-judge agreement can
be read against an attenuation ceiling rather than against $1.0$, which is what makes the
one rubric that genuinely disagrees (MITI) distinguishable from the ones that are merely
noisy.
